% GENERATED FILE. Edit canonical lessons, then run npm run generate:books.
% canonical-source-hash: fnv1a64:fa4a1931f9457c57
% canonical-lessons: KA-C484-tapamana, KA-C484-ardrate, KA-C484-havaguna, KA-C484-teva, KA-C484-hadagedu

\chapter{The Weather Report}
\label{ch:ka-the-weather-report}
\hlchaptermodality{484}

\begin{chapteropening}
\textbf{By the end of this chapter:} \emph{I can say temperature, humidity, climate, dampness and to worsen.}

Complete the last lesson of chapter 484: I can say temperature, humidity, climate, dampness and to worsen.
\end{chapteropening}

\section[\texorpdfstring{\kn{ತಾಪಮಾನ} (tāpamāna)}{ತಾಪಮಾನ (tāpamāna)}]{\kn{ತಾಪಮಾನ} (tāpamāna) — temperature}

\label{lesson:KA-C484-tapamana}



\begin{quote}
Before the new one: say the Kannada for groundwater, then the Kannada for to evaporate.
\end{quote}

\subsection*{You'll want to know: \kn{ತಾಪಮಾನ}}
\textbf{\kn{ತಾಪಮಾನ}} — \emph{tāpamāna} — \textquotedblleft{}temperature\textquotedblright{}.

\begin{quote}
\kn{ಇವತ್ತು} \kn{ತಾಪಮಾನ} \kn{ನಿನ್ನೆಗಿಂತ} \kn{ಹೆಚ್ಚು}.
\end{quote}

\begin{quote}
\emph{ivattu tāpamāna ninneginta heccu.}
\end{quote}

\begin{quote}
Today the temperature is higher than yesterday.
\end{quote}

\kn{ತಾಪ}, heat, plus \kn{ಮಾನ}, a measure.

\begin{cousinweb}
\textbf{\kn{ತಾಪಮಾನ}} (\emph{tāpamāna}) is a modern compound built from Sanskrit words: \emph{tāpa}, \textquotedblleft{}heat, glow\textquotedblright{} (CDIAL 5767), and \emph{māna}, \textquotedblleft{}a measure, a standard\textquotedblright{} (CDIAL 10041).
\end{cousinweb}

\begin{grammarlens}[title={What you've built}]
One more word for getting around.
\end{grammarlens}

\subsection*{Your turn}
\begin{itemize}
\raggedright
  \item \emph{Say it:} \emph{tāpamāna}
  \item \emph{Say it:} \emph{tāpamāna}, once more
  \item \emph{Say it:} \emph{ivattu tāpamāna ninneginta heccu.} — Today the temperature is higher than yesterday.
  \item \emph{Recall:} say the Kannada for groundwater, then the Kannada for to evaporate, then say \emph{tāpamāna} again
\end{itemize}

\begin{tcolorbox}[breakable,colback=teal!4,colframe=teal!35!black,title={Before you move on}]
What does \kn{ತಾಪಮಾನ} mean? (\textquotedblleft{}Temperature\textquotedblright{}.) Say it once more.
\end{tcolorbox}
\section[\texorpdfstring{\kn{ಆರ್ದ್ರತೆ} (ārdrate)}{ಆರ್ದ್ರತೆ (ārdrate)}]{\kn{ಆರ್ದ್ರತೆ} (ārdrate) — humidity}

\label{lesson:KA-C484-ardrate}



\begin{quote}
Before the new one: say the Kannada for to evaporate, then the Kannada for temperature.
\end{quote}

\subsection*{You'll want to know: \kn{ಆರ್ದ್ರತೆ}}
\textbf{\kn{ಆರ್ದ್ರತೆ}} — \emph{ārdrate} — \textquotedblleft{}humidity\textquotedblright{}.

\begin{quote}
\kn{ಕರಾವಳಿಯಲ್ಲಿ} \kn{ಆರ್ದ್ರತೆ} \kn{ಹೆಚ್ಚು}, \kn{ಅದಕ್ಕೇ} \kn{ಸೆಕೆ} \kn{ಆಗುತ್ತದೆ}.
\end{quote}

\begin{quote}
\emph{karāvaḷiyalli ārdrate heccu, adakkē seke āguttade.}
\end{quote}

\begin{quote}
Humidity is high on the coast, and that is why it feels sticky.
\end{quote}

Weather reports give it in per cent.

\begin{cousinweb}
\textbf{\kn{ಆರ್ದ್ರತೆ}} (\emph{ārdrate}) is built from Sanskrit \emph{ārdra}, \textquotedblleft{}wet\textquotedblright{} (CDIAL 1340), with the noun ending -\emph{tā}, \textquotedblleft{}-ness\textquotedblright{}: wetness.
\end{cousinweb}

\begin{grammarlens}[title={What you've built}]
One more word for getting around.
\end{grammarlens}

\subsection*{Your turn}
\begin{itemize}
\raggedright
  \item \emph{Say it:} \emph{ārdrate}
  \item \emph{Say it:} \emph{ārdrate}, once more
  \item \emph{Say it:} \emph{karāvaḷiyalli ārdrate heccu, adakkē seke āguttade.} — Humidity is high on the coast, and that is why it feels sticky.
  \item \emph{Recall:} say the Kannada for to evaporate, then the Kannada for temperature, then say \emph{ārdrate} again
\end{itemize}

\begin{tcolorbox}[breakable,colback=teal!4,colframe=teal!35!black,title={Before you move on}]
What does \kn{ಆರ್ದ್ರತೆ} mean? (\textquotedblleft{}Humidity\textquotedblright{}.) Say it once more.
\end{tcolorbox}
\section[\texorpdfstring{\kn{ಹವಾಗುಣ} (havāguṇa)}{ಹವಾಗುಣ (havāguṇa)}]{\kn{ಹವಾಗುಣ} (havāguṇa) — climate}

\label{lesson:KA-C484-havaguna}



\begin{quote}
Before the new one: say the Kannada for temperature, then the Kannada for humidity.
\end{quote}

\subsection*{You'll want to know: \kn{ಹವಾಗುಣ}}
\textbf{\kn{ಹವಾಗುಣ}} — \emph{havāguṇa} — \textquotedblleft{}climate\textquotedblright{}.

\begin{quote}
\kn{ಗಿರಿಧಾಮದ} \kn{ಹವಾಗುಣ} \kn{ತಣ್ಣಗಿದೆ}.
\end{quote}

\begin{quote}
\emph{giridhāmada havāguṇa taṇṇagide.}
\end{quote}

\begin{quote}
The climate of the hill station is cool.
\end{quote}

\kn{ಹವಾಮಾನ} is the weather of a day; \kn{ಹವಾಗುಣ} is the usual weather of a place.

\begin{cousinweb}
\textbf{\kn{ಹವಾಗುಣ}} (\emph{havāguṇa}) is a Kannada compound of \kn{ಹವಾ}, \textquotedblleft{}air, weather\textquotedblright{}, usually traced to Arabic \emph{hawā}, and Sanskrit \emph{guṇa}, \textquotedblleft{}a quality\textquotedblright{} (CDIAL 4190).
\end{cousinweb}

\begin{grammarlens}[title={What you've built}]
One more word for getting around.
\end{grammarlens}

\subsection*{Your turn}
\begin{itemize}
\raggedright
  \item \emph{Say it:} \emph{havāguṇa}
  \item \emph{Say it:} \emph{havāguṇa}, once more
  \item \emph{Say it:} \emph{giridhāmada havāguṇa taṇṇagide.} — The climate of the hill station is cool.
  \item \emph{Recall:} say the Kannada for temperature, then the Kannada for humidity, then say \emph{havāguṇa} again
\end{itemize}

\begin{tcolorbox}[breakable,colback=teal!4,colframe=teal!35!black,title={Before you move on}]
What does \kn{ಹವಾಗುಣ} mean? (\textquotedblleft{}Climate\textquotedblright{}.) Say it once more.
\end{tcolorbox}
\section[\texorpdfstring{\kn{ತೇವ} (tēva)}{ತೇವ (tēva)}]{\kn{ತೇವ} (tēva) — dampness; moisture}

\label{lesson:KA-C484-teva}



\begin{quote}
Before the new one: say the Kannada for humidity, then the Kannada for climate.
\end{quote}

\subsection*{You'll want to know: \kn{ತೇವ}}
\textbf{\kn{ತೇವ}} — \emph{tēva} — \textquotedblleft{}dampness; moisture\textquotedblright{}.

\begin{quote}
\kn{ಮಳೆಗಾಲದಲ್ಲಿ} \kn{ಗೋಡೆಗಳಲ್ಲಿ} \kn{ತೇವ} \kn{ಇರುತ್ತದೆ}.
\end{quote}

\begin{quote}
\emph{maḷegāladalli gōḍegaḷalli tēva iruttade.}
\end{quote}

\begin{quote}
In the rainy season there is damp in the walls.
\end{quote}

\kn{ತೇವದ} \kn{ಬಟ್ಟೆ}, damp clothes.

\begin{cousinweb}
\textbf{\kn{ತೇವ}} (\emph{tēva}) is a Kannada word for dampness whose origin is not settled.
\end{cousinweb}

\begin{grammarlens}[title={What you've built}]
One more word for getting around.
\end{grammarlens}

\subsection*{Your turn}
\begin{itemize}
\raggedright
  \item \emph{Say it:} \emph{tēva}
  \item \emph{Say it:} \emph{tēva}, once more
  \item \emph{Say it:} \emph{maḷegāladalli gōḍegaḷalli tēva iruttade.} — In the rainy season there is damp in the walls.
  \item \emph{Recall:} say the Kannada for humidity, then the Kannada for climate, then say \emph{tēva} again
\end{itemize}

\begin{tcolorbox}[breakable,colback=teal!4,colframe=teal!35!black,title={Before you move on}]
What does \kn{ತೇವ} mean? (\textquotedblleft{}Dampness; moisture\textquotedblright{}.) Say it once more.
\end{tcolorbox}
\section[\texorpdfstring{\kn{ಹದಗೆಡು} (hadageḍu)}{ಹದಗೆಡು (hadageḍu)}]{\kn{ಹದಗೆಡು} (hadageḍu) — to worsen; to deteriorate}

\label{lesson:KA-C484-hadagedu}



\begin{quote}
Before the new one: say the Kannada for climate, then the Kannada for dampness.
\end{quote}

\subsection*{You'll want to know: \kn{ಹದಗೆಡು}}
\textbf{\kn{ಹದಗೆಡು}} — \emph{hadageḍu} — \textquotedblleft{}to worsen; to deteriorate\textquotedblright{}.

\begin{quote}
\kn{ಸಂಜೆಯ} \kn{ಹೊತ್ತಿಗೆ} \kn{ಹವಾಮಾನ} \kn{ಹದಗೆಟ್ಟಿತು}.
\end{quote}

\begin{quote}
\emph{sañjeya hottige havāmāna hadageṭṭitu.}
\end{quote}

\begin{quote}
By evening the weather had turned bad.
\end{quote}

Past (it …): \kn{ಹದಗೆಟ್ಟಿತು}, \emph{hadageṭṭitu}.

Of weather, roads and health.

\begin{cousinweb}
\textbf{\kn{ಹದಗೆಡು}} (\emph{hadageḍu}) is native Dravidian: \emph{hada}, \textquotedblleft{}the proper state\textquotedblright{} (DEDR 3907 lists Kannada \emph{pada}, \emph{hadanu}), and \emph{keḍu}, \textquotedblleft{}to be ruined, to spoil\textquotedblright{} (DEDR 1942): to fall out of its proper state.
\end{cousinweb}

\begin{grammarlens}[title={What you've built}]
One more word for getting around.
\end{grammarlens}

\subsection*{Your turn}
\begin{itemize}
\raggedright
  \item \emph{Say it:} \emph{hadageḍu}
  \item \emph{Say it:} \emph{hadageḍu}, once more
  \item \emph{Say it:} \emph{sañjeya hottige havāmāna hadageṭṭitu.} — By evening the weather had turned bad.
  \item \emph{Recall:} say the Kannada for climate, then the Kannada for dampness, then say \emph{hadageḍu} again
\end{itemize}

\begin{tcolorbox}[breakable,colback=teal!4,colframe=teal!35!black,title={Before you move on}]
What does \kn{ಹದಗೆಡು} mean? (\textquotedblleft{}To worsen; to deteriorate\textquotedblright{}.) Say it once more.
\end{tcolorbox}
